% \section{Equity and Affirmative Action}
\section{Defining Equity}
Equity is the distribution of resources among groups to overcome obstacles and to raise their opportunities for access \cite{schement2001imagining}. Thus, historically disadvantaged groups are compensated, and others get their fair share to reach their goals. To operationalize equity, we learn the existing biases from the historical data and compensate for them in the predictions generated by the model. This can be viewed similarly to affirmative action in which present decisions (algorithmic or otherwise) are made to compensate for biases of the past but in such a way that groups reach an ultimate equality. The goal is to equalize all the groups in the long run and give each group their fair share.% rather than equalizing everyone blindly which can result in some groups getting excessive and unfair attention over others. %% FM: I don't understnad what you're trying to say here so i deleted it
We modeled this phenomenon in our classification objective function which will be described in detail in this section. \\ 
Our objective function consists of two terms:
\begin{itemize}
    \item The Fairness Objective: In which the goal is to enforce equity amongst groups.
    \item The Classification Objective: In which we enforce the classification objective to achieve predictive accuracy.
\end{itemize}
Finally, we combine these two objectives and control the importance of each using a regularization parameter. % $\beta$.
\subsection{Operationalizing Equity}
Herein, we formalize our equity notion of fairness. Let $Y$ be a random variable denoting an outcome of a decision making process and $A$ be the sensitive variable (e.g., demographic group membership).  
Let $D$ be the set of all decisions made in the past (or historical decisions).
Let the joint distribution $p_D(Y, A)$ summarize the essential statistics of past decisions.
We are interested in the case when past decisions contain bias, i.e., $p_D(Y\mid A=a) \not\approx p_D(Y \mid A=b)$.

Let $M$ be the set of instances for which a prediction has to be done using a machine learning method.
As mentioned earlier we want the decisions of the classifier to account for and reverse the biases in the data in an equitable manner.
We formalize this in the following way.
Let a joint distribution $p_M(Y, A)$ summarize the essential statistics of the classifier on $M$.
Our goal is to generate equitable decisions for each group:
\begin{align*}
&p_D(Y=y \mid A=a) p_D(A=a) + p_M(Y = y \mid A=a) p_M(A=a)\\ 
=&p_D(Y=y \mid A=b) p_D(A=b) + p_M(Y = y \mid A=b) p_M(A=b), 
\end{align*}

for each possible outcome $y$.
We assume that we study the same number of instances from each demographic group between both the historical and outcome sets (i.e., $p_D(A = a) = p_D(A = b) = 1/2$ and $p_M(A=a)=p_M(A=b)=1/2$).
Under this assumption, our fairness criterion becomes:
\begin{align*}
&p_D(Y=y \mid A=a) + p_M(Y = y \mid A=a)\\ 
=&p_D(Y=y \mid A=b)  + p_M(Y = y \mid A=b). 
\end{align*}
One can interpret this criterion as equalizing odds of different groups for a random person drawn with equal probability from $D$ or $M$.

To comply with the notation widely used in the fairness literature,  $p_D(Y | A=a)$ can also be written as $p(Y|A=a)$ and $p_M(Y| A=a)$ as $p(\hat{Y}| A=a)$.\\
Finally, to translate our new fairness notion into widely used format one can write our objective as follows:
\begin{align*}
    p(\hat{Y}|A=a)+p(Y|A=a)=p(\hat{Y}|A=b)+p(Y|A=b).
\end{align*}
\begin{tcolorbox}
\begin{definition} [Statistical Equity]
  A predictor is statically equitable among demographic groups, $a$ and $b$, if it satisfies $p(\hat{Y}|A=a)+p(Y|A=a)=p(\hat{Y}|A=b)+p(Y|A=b)$.
\end{definition}
%\todo{NEW NOTATIONS ENDS HERE}
\end{tcolorbox}




% Let $O$ denote the set of outcomes in a decision making process, $A$ denote the demographic or group membership information, $D$ denote the distribution that the data is coming from i.e. $D=h$ denotes the historical distribution in the train data and $D=m$ denotes the future or outcome distribution in the test data whose labels would be based upon predictive outcomes of the model. \\
% Notice that due to the feedback loop property $M$ can be part of $H$, but initially they are disjoint. To comply with the notation widely used in the fairness literature,  $p(o \in O| A=a,D=h)$ can also be written as $p(Y|A=a)$ and $p(o \in O| A=a,D=m)$ as $p(\hat{Y}| A=a)$. However, for our preliminary analysis, we will utilize our current notation and later translate our notation into fairness language for ease of following.\\
% Note: We use parity and equality interchangeably as synonyms throughout this paper.
% \subsection{Fairness Objective}
% In order to model the fairness objective to consider and correct the historical biases in future decision making outcomes, we marginalize the outcome distribution over the two possible distributions, historical and future (h and m).
%   \begin{align*}
%     p(o \in O|A=a) &= p(o \in O| A=a,D=h)p(D=h|A=a) \\
%   &+ p(o \in O| A=a,D=m)p(D=m|A=a)
% \end{align*}
% Thus to satisfy fairness amongst different demographic groups a and b, we would have the following:
%  \begin{align*}
%     \sum_{d=h,m}p(o \in O| A=a,D=d)p(D=d|A=a) \\
%     =\sum_{d=h,m}p(o \in O| A=b,D=d)p(D=d|A=b)
% \end{align*}
% To relax the problem, we assume that we study the same number of instances from each demographic group in both historical and outcome distributions. This leads us to the relaxed version of our fairness objective as follows:
% \begin{align*}
%      p(o \in O| A=a,D=h)+ p(o \in O| A=a,D=m)\\ 
%      =p(o \in O| A=b,D=h)+ p(o \in O| A=b,D=m)
% \end{align*}
% Finally, to translate our new fairness notion into widely used format one can write our objective as follows:
% \begin{align*}
%     p(\hat{Y}|A=a)+p(Y|A=a)=p(\hat{Y}|A=b)+p(Y|A=b)
% \end{align*}


\subsection{Classification Objective}
In order to satisfy the fairness objective %among demographic groups a and b: $p(\hat{Y}|A=a)+p(Y|A=a)=p(\hat{Y}|A=b)+p(Y|A=b)$ as defined earlier, 
and to be able to couple it with our classification loss, we divided our total classification objective into two terms. The first term, denoted by $F_{equity}$, attempts to satisfy the fairness objective, and the other term, denoted by $L$, attempts to satisfy the classification loss. These two losses are then controlled and coupled by a regularization term $\beta$.

The resulting fairness objective for each possible outcome $y$ is coupled in our classifier as follows:
%  \begin{align*}
%         F_\text{equity}(\theta)=&\sum_{y}\left(\left[\frac{1}{n}\sum_{i=1}^{n}p(\hat{Y_i}=y|A=a)
%         +p(Y=y|A=a)\right]\right. \\
%         &\left. -\left[\frac{1}{n}\sum_{i=1}^{n}p(\hat{Y_i}=y|A=b) 
%         +p(Y=y|A=b)\right]\right)^2. \\
% \end{align*}
 \begin{align*}
        &F_\text{equity}(\theta)=\\
        &\;\;\;\;\;\;\sum_{y}\left(\Big[\frac{1}{n}\sum_{i=1}^{n}p(\hat{Y_i}=y|A=a)
        +p(Y=y|A=a)\Big]\right. \\
        &\;\;\;\;\;\;\left. -\Big[\frac{1}{n}\sum_{i=1}^{n}p(\hat{Y_i}=y|A=b) 
        +p(Y=y|A=b)\Big]\right)^2.
\end{align*}

Note that there is a difference between the notation of historical and predictive (future) outcomes. By equalizing the sum of historical plus future outcomes of one group to another, we are enforcing affirmative action and try to compensate for observed historical biases in the data by correcting and adjusting the predictive outcome so that eventually all the groups reach an equilibrium in our objective function. This equilibrium can be in terms of all the groups satisfying their goals, e.g. that in Figure \ref{motivation}. 
We defined our classification objective as the cross-entropy loss $L(\theta)$.
%\begin{equation}
    %\begin{multlined}
    % \begin{align*}
    %     L(\theta) = -Y\log(\hat{Y})-(1-Y)\log(1-\hat{Y}).
    % \end{align*}
    % \begin{align*}
    %     L(\theta) =  - \sum_{i=1}^n \log p_\theta(\widehat{Y}_i = Y_i \mid X_i).
    % \end{align*}
    %\end{multlined}
%\end{equation}
Notice that $L(\theta)$ can be any other loss; however, cross-entropy is used in our experiments. 
Finally, by combining the fairness objective with the classification objective, we define the whole objective of our fair classification task as follows: \\
\begin{equation}
     \min_\theta \; \beta F_\text{equity}(\theta)+ (1-\beta)L(\theta).
     \label{equity_eq}
\end{equation}
$\beta$ is a hyperparameter that controls the importance of the fairness constraint over the classification objective. By making the $\beta$ value larger we are enforcing more of the fairness (equity) constraint in our objective, while smaller $\beta$ value favors the classification accuracy over the fairness constraint.
